% ============================================================
% Section 5: Conclusion
% ============================================================
\section{Conclusion}

\begin{frame}{Summary}
    \textbf{Mean Shift Clustering:}
    \begin{itemize}
        \item \textcolor{customred}{Core idea:} Find density modes, points converge to nearest mode
        \item \textcolor{customred}{Strengths:} No cluster count needed, flexible shapes, robust
        \item \textcolor{customred}{Weaknesses:} $O(n^2)$ complexity, bandwidth sensitivity, high-dim issues
    \end{itemize}

    \vspace{0.5cm}

    % \textbf{Experimental Findings:}
    % \begin{itemize}
    %     \item Outperforms K-means in cluster quality (synthetic \& Iris)
    %     \item Requires dimensionality reduction for high-dim data (Wine)
    %     \item Excellent for image segmentation, but slower
    %     \item Bandwidth selection is critical
    % \end{itemize}

    % \vspace{0.5cm}

    \textbf{Best Use Cases:}
    \begin{itemize}
        \item Small-medium datasets with unknown cluster count
        \item Non-spherical clusters or mode-finding tasks
        \item Image/video processing applications
    \end{itemize}
\end{frame}

\begin{frame}{References}
    \begin{enumerate}
        \item Fukunaga, K., \& Hostetler, L. (1975). \textit{The estimation of the gradient of a density function, with applications in pattern recognition}. IEEE Transactions on Information Theory.

        \item Comaniciu, D., \& Meer, P. (2002). \textit{Mean shift: A robust approach toward feature space analysis}. IEEE Transactions on Pattern Analysis and Machine Intelligence.

    \end{enumerate}
\end{frame}

\begin{frame}[plain]
    \begin{center}
        \vspace{2cm}

        {\Huge \textcolor{customred}{Thank You!}}

        \vspace{1cm}

        {\Large Questions?}

        \vspace{1.5cm}

        {\normalsize
        Group K + F \\
        Computational Statistics Lab
        }
    \end{center}
\end{frame}
